\clearpage
\fancyhead[L]{\small EFFEMERIDI ASTRONOMICHE DI MILANO PER L'ANNO 1817.}
\hypertarget{native-606-lower}{}
\begin{center}\small Göttingische gelehrte Anzeigen. Stück 144. Seiten 1433..1438. 8. September 1817.\\[3pt]\rule{.72\linewidth}{.3pt}\end{center}

{\setlength{\rightskip}{0pt plus 3em}\emph{Effemeridi astronomiche di Milano per l'anno \textup{1817} calcolate da} \textsc{Francesco Carlini} \emph{ed} \textsc{Enrico Brambilla}. \emph{Con Appendice.} Aus der K. K. Druckerei. Mailand 1816. 108 und 116 Seiten im Oktavformat.\par}

\begingroup\brokenpenalty=10000\relax
Da die Einrichtung des Kalenders in diesem Jahrgang und die beigefügten Zusätze ganz dieselben sind wie beim vorhergehenden, so haben wir diesmal bloß die Aufsätze des Anhangs anzuzeigen. Im ersten liefert uns Herr \textsc{Oriani} das Verzeichnis der Deklinationen von 40 Sternen aus den in früheren Jahrgängen der Ephemeriden abgedruckten Beobachtungen mit dem drei Fuß großen Vervielfältigungskreis von \textsc{Reichenbach}. Der beigefügte Vergleich mit den Angaben von \textsc{Maskelyne}, \textsc{Piazzi} und \textsc{Pond} ist wegen der von diesen Astronomen angewandten Instrumente bemerkenswert. In Hinsicht auf die erreichbare Genauigkeit setzt Hr.\ \textsc{Oriani} den drei Fuß großen Wiederholungskreis von \textsc{Reichenbach} dem fünf Fuß großen Mauerkreis von \textsc{Troughton} in Greenwich ungefähr gleich, und in Hinsicht auf die Bequemlichkeit des Beobachtens räumt er gleichfalls dem letzteren Instrument keinen Vorzug ein, indem das, was ein viermaliges Wiederholen an Zeit mehr kostet, durch das am Wiederholungskreis bequemere Ablesen ausgeglichen werde. Hierbei ist freilich der zur Berechnung der außerhalb des Meridians gemachten Beobachtungen erforderliche Zeitaufwand nicht in Anschlag gebracht. --- Der zweite Aufsatz, von \textsc{Ottaviano Fabrizio Mossotti}, neue Analyse der Aufgabe, die Bahnen der Himmelskörper zu bestimmen, verdient wegen der Wichtigkeit des Gegenstandes eine etwas ausführlichere Darstellung. Die ganze Behandlung ist rein analytisch und mit Geschicklichkeit durchgeführt. Indessen wird durch die große Menge der Zeichen die Festhaltung des Fadens, die Übersicht des Ganzen und die klare Auffassung dessen, worauf es eigentlich ankommt, etwas erschwert. Das Wesentliche beruht auf Folgendem. Wenn die Hälften der Bahnparameter der zu bestimmenden Bahn und der Erdbahn mit $p$ und $P$, die Produkte von $\BareRoot{p}$ \SourceLeaf{607}mit den Kosinus der Neigung der Ebene der unbekannten Bahn gegen drei feste willkürliche Ebenen mit $c$, $c'$, $c''$, und ebenso die Produkte von $\BareRoot{P}$ mit den Kosinus der Neigungen der Erdbahn gegen dieselben Fundamentalebenen mit $C$, $C'$, $C''$ bezeichnet werden, so findet Herr \textsc{Mossotti} drei lineare Gleichungen von der Form
\[
A(c-C)+A'(c'-C')+A''(c''-C'') = 0
\]
wo die Koeffizienten bekannte Größen sind. Diese Gleichungen sind jedoch nur näherungsweise richtig, insofern die Zwischenzeiten nicht unendlich klein sind, und gründen sich eigentlich bloß darauf, dass die drei Flächensegmente, die entstehen, wenn die drei Orte paarweise durch Sehnen verbunden werden, sowohl bei der unbekannten Bahn als auch bei der Erdbahn sehr nahe den Kuben der Zwischenzeiten proportional sind. Eben deswegen, weil jene Gleichungen nur näherungsweise wahr sind, ist es nicht gestattet, sie so zu kombinieren, dass eine der Größen $c-C$, $c'-C'$, $c''-C''$ eliminiert würde, sondern es lässt sich zeigen, dass alle drei Gleichungen eigentlich nur für eine gelten dürfen, und Herr \textsc{Massotti} behält daher auch nur eine von ihnen bei. Da nun diese für sich allein nicht weiterführen kann, nimmt Herr \textsc{Massotti} eine vierte Beobachtung zu Hilfe, die, mit zweien der vorigen verbunden, eine ganz ähnliche Gleichung liefert, deren Verbindung mit der aus den drei ersten Beobachtungen gefolgerten allerdings rechtmäßig ist und zur Bestimmung des Verhältnisses der drei Größen $c-C$, $c'-C'$, $c''-C''$ untereinander benutzt wird, sodass die Resultate in dieser Gestalt erscheinen
\[
c'-C' = M(c-C), \quad c''-C'' = N(c-C)
\]
Durch Benutzung des oben ausgesprochenen Prinzips findet hernach Herr \textsc{Massotti} auch das Verhältnis der drei Abstände des Himmelskörpers von der Erde zu einer der drei Größen $c-C$, $c'-C'$, $c''-C''$, und endlich die absoluten Werte dieser drei Größen selbst. Da nun $C$, $C'$, $C''$ für sich bekannt sind, ergeben sich die Werte von $c$, $c'$, $c''$, und damit $p = cc+c'c'+c''c''$, sowie aus
\[
\frac{c}{\BareRoot{p}}, \quad \frac{c'}{\BareRoot{p}}, \quad \frac{c''}{\BareRoot{p}}
\]
die Größen, von denen die Lage der Bahn des Himmelskörpers abhängt. Die weitere Bestimmung der übrigen Elemente beruht auf bekannten Sätzen und hat nichts Eigentümliches.

Sollen wir nun unser Urteil über diese neue Methode hier abgeben, so erkennen wir zunächst mit Vergnügen an, dass die erwähnte lineare Gleichung zwischen $c-C$, $c'-C'$, $c''-C''$ bemerkenswert, und dass es interessant ist, die Möglichkeit einer durchgehend bloß linearen Bestimmung einer unbekannten Bahn entwickelt zu sehen. Eine andere Frage ist nun aber, ob eine solche Methode zur wirklichen praktischen Anwendung zu empfehlen sei. Es ist klar, dass hier nicht alles benutzt ist, und nicht benutzt werden sollte, was in der Theorie der Bewegung in Kegelschnitten liegt, und zwar wird namentlich, um es im Geiste einer Näherungsmethode auszusprechen, die Bedingung ganz ignoriert, dass die Quotienten, wenn jene Segmente mit den Würfeln der Zwischenzeiten dividiert werden, sehr nahe gleich sein müssen der Größe $\dfrac{k^3}{12r^3}$, wo $k$ in der Bedeutung der Theoria Motus Corporum Coelestium zu nehmen ist, und $r$ den Abstand des Himmelskörpers von der Sonne bedeutet. Daher also die Notwendigkeit, eine vierte Beobachtung zu Hilfe zu nehmen bei einer Aufgabe, wo drei Beobachtungen schon die vollständige Auflösung einschließen. Auch abgesehen davon, dass dies in theoretischer Rücksicht nicht angemessen ist, wird dieses Verfahren in praktischer Hinsicht deswegen sehr misslich, weil eigentlich hier die Bestimmung der Bahn auf Größen der \emph{dritten} Ordnung beruht, und daher die unvermeidlichen Beobachtungsfehler einen unverhältnismäßigen Einfluss auf die Resultate äußern müssen, der desto größer sein wird, weil die Methode, als bloße Näherungsmethode, nur kurze \SourceLeaf{608}Zwischenzeiten anzuwenden gestatten kann. Diese Methode wird daher im Allgemeinen nur wenig zuverlässige Resultate geben können, wo die vollständige Benutzung dreier Beobachtungen schon zu einer sehr genauen Näherungsbestimmung führen kann.

Endlich müssen wir noch bemerken, dass, wenn man einmal bloß von jenem Prinzip der Proportionalität der Segmente zu den Würfeln der Zwischenzeiten ausgehen will, man denselben Erfolg einer bloß linearen Bestimmung der Bahn aus vier Beobachtungen auf eine einfachere Art haben kann. Jenes Prinzip gibt nämlich aus drei Beobachtungen mit leichter Mühe eine gleichfalls lineare Gleichung zwischen zwei Abständen von der Erde, und zwischen denselben Abständen findet man auf dieselbe Weise eine zweite ähnliche Gleichung, indem man die zugehörigen beiden Beobachtungen mit einer vierten Beobachtung verbindet. Die Elimination gibt dann diese Abstände selbst, ebenso genau wie die Werte, welche dieselben bei der Anwendung von \textsc{Mossotti}s Methode erhalten, oder vielmehr eigentlich dieselben Werte. Sobald diese Abstände bekannt sind, hat die Bestimmung der ganzen Bahn bekanntlich keine Schwierigkeit. Die weitere Entwicklung dieses Gegenstandes, wozu natürlich hier nicht der Ort ist, müssen wir auf eine andere Gelegenheit versparen. Nur das eine müssen wir noch bemerken, dass der im vorliegenden Jahrgang der Ephemeriden abgedruckte 80 Seiten lange Aufsatz nur das Theoretische von \textsc{Mossotti}s Methode enthält, und also eine versprochene Anwendung auf den \textsc{Halley}’schen Kometen wahrscheinlich im nächsten Jahrgang nachfolgen wird. --- Den Schluss des vorliegenden Jahrgangs macht ein kleiner Artikel von \textsc{Carlini}, worin einige Fehler in \textsc{Delambre}s Tafeln für die Jupitersatelliten berichtigt werden.
\par\clearpage\endgroup
